\documentclass[11pt]{article}
\usepackage{amsmath,amssymb,amsthm,booktabs,geometry,hyperref}
\geometry{margin=.88in}\hypersetup{hidelinks}
\newtheorem{theorem}{Theorem}\newtheorem{lemma}{Lemma}\newtheorem{corollary}{Corollary}
\title{The sharp fractional plank bound for packed triangles in three coordinates}
\author{Leo Gao\\\small AI-assisted research note}\date{}
\begin{document}\maketitle
\begin{abstract}
Aharoni, Holzman, Krivelevich and Meshulam conjectured that any fractional plank cover of a convex set meeting every facet of the unit $n$-cube costs at least $2/n$. We prove the sharp $2/3$ bound for every triangle meeting all six facets of the three-cube, without a symmetry assumption. For triangles with three distinct coordinate-median vertices we construct an inscribed triangle with simultaneous uniform marginals of common width at least $2/3$; all other triangles contain a cube vertex and are treated by a different argument. We also give a gluing criterion for full-dimensional cyclically invariant sets. No general result on all planar sets, all bodies in three dimensions, or ordinary plank covers of triangles is claimed.
\end{abstract}
A set in $[0,1]^n$ is \emph{packed} if it meets the $2n$ facets. A fractional plank cover of $K$ is a family of nonnegative measurable $f_j$ on $[0,1]$ satisfying $\sum f_j(x_j)\ge1$ throughout $K$; its cost is $\sum\int_0^1 f_j$. The conjecture is from \cite[Conj.~3.4]{AHKM}, and is recorded in \cite[Sec.~3.4]{BMP}. It concerns arbitrary measurable functions, unlike the ordinary three-strip question.
\begin{theorem}\label{main} Every fractional plank cover of a packed triangle in $[0,1]^3$ has cost at least $2/3$. The bound is sharp.\end{theorem}
\begin{lemma}\label{distinct} Suppose the three vertices assigned to the median in each respective coordinate can be chosen distinct. There exists a probability on a subtriangle having uniform marginals of common width at least $2/3$. The width is strictly greater than $2/3$ if no coordinate has a tie.\end{lemma}
\begin{proof} Label the median vertices $1,2,3$ and reflect coordinates so the vertices are $(a,0,1),(1,b,0),(0,1,c)$, respectively. We work first with $0<a,b,c<1$. Let $L$ send the three standard barycentric vertices to these physical vertices, extended linearly on $\mathbb R^3$. Put
\[
 A=1-b+bc,\quad B=1-c+ca,\quad C=1-a+ab.
\]
Use tangent directions on the barycentric simplex
\[
 d_1=(1,-a,a-1),\quad d_2=(b-1,1,-b),\quad d_3=(-c,c-1,1).
\]
The columns $Ld_i$ form the matrix
\[
 \begin{pmatrix}0&C&-B\\-C&0&A\\B&-A&0\end{pmatrix}.
\]
Its null vector $(A,B,C)$ is strictly positive. The same relation holds among $d_i$, since the displayed matrix has rank two and $L$ is therefore injective on their two-dimensional span. Put
\[
 u=\frac{C+B-A}{2CB},\quad v=\frac{A+C-B}{2AC},\quad w=\frac{A+B-C}{2AB}.
\]
The following exact polynomial inequalities are verified below:
\begin{equation}\label{key}
 u,v,w>0,\qquad M:=\max((1-b)v,cw)+\max(au,(1-c)w)+\max((1-a)u,bv)<3/2.
\end{equation}
Let $t=1/M$ and let $m$ have as barycentric coordinates $t$ times the three respective maxima in \eqref{key}. They sum to one. The three points $s_1=m+tud_1$, $s_2=m+tvd_2$, $s_3=m+twd_3$ have nonnegative coordinates by definition of the maxima and belong to the barycentric simplex. The positive dependence shows $m\in\operatorname{conv}\{s_i\}$. In physical coordinate $j$, $L(s_j)_j=L(m)_j$ and the other vertices are strictly on opposite sides. All three projection widths equal $t$ because
\[
 Cv+Bw=Cu+Aw=Bu+Av=1.
\]

Take the barycentric coordinates $q_i$ of $m$ in this new triangle, and assign mass $q_i$ uniformly to the edge opposite $L(s_i)$. In one normalized projection with vertex values $0,h,1$, denote the weights by $q_0,q_h,q_1$. We have $q_1+hq_h=h$, since the barycenter is at the median level. The densities on $(0,h)$ and $(h,1)$ are $q_h+q_1/h$ and $q_h+q_0/(1-h)$, respectively; both equal one. Each physical marginal is thus uniform of width $t>2/3$.

For boundary parameters approximate by interior triples. By compactness extract a subsequence on which all vertices, weights and interval endpoints converge. A uniform edge law is the affine image of a single uniform parameter, hence converges weakly even if the edge collapses. Its mixture converges on the limiting subtriangle, and the coordinate laws remain uniform on limiting intervals of common width at least $2/3$.
\end{proof}
\paragraph{Polynomial verification of \eqref{key}.} Let $D=2ABC$; set $(U,V,W)=D(u,v,w)$. Check positivity of $U,V,W$ and of each of the eight polynomials
\[
 3ABC-e_1-e_2-e_3,\quad
 (e_1,e_2,e_3)\in\{(1-b)V,cW\}\times\{aU,(1-c)W\}\times\{(1-a)U,bV\}.
\]
The following integer coefficient certificate is also executable as \texttt{distinct\_verify.py}. For $P=\sum p_{ijk}a^ib^jc^k$ of coordinate degrees $(d,e,f)$ collect the coefficients of
\[
 \sum p_{ijk}x^iy^jz^k(1+x)^{d-i}(1+y)^{e-j}(1+z)^{f-k}.
\]
For the three numerators followed by the eight differences in lexicographic binary choice order, the counts of nonzero monomials are
\[
10,10,10;\quad 17,17,14,17,17,14,17,17.
\]
In every case the smallest nonzero coefficient is $1$. The displayed pullback is therefore positive for positive $x,y,z$. Substituting $x=a/(1-a)$ and cyclically gives the assertion. No numerical optimizer is used.
\begin{lemma}\label{corner} If a packed $K\subseteq[0,1]^3$ is convex and contains a cube vertex, its fractional plank covering number is at least $2/3$.\end{lemma}
\begin{proof} Reflect so the vertex is $0$ and choose $p_i\in K$ with $(p_i)_i=1$. We first show no ordinary measurable coordinate-plank cover can have total width $<1$. Otherwise, enlarge the forbidden sets on the real line to open sets still of total length $<1$. Their components are intervals; take a finite subcover of the compact hull of $0,p_1,p_2,p_3$. Write the selected intervals as $(a_k,b_k)$ in coordinates $i(k)$. Choose $L_k>b_k-a_k$ with $\sum L_k<1$.

For a binary vector $s$ put $x(s)=\sum_k L_k s_k p_{i(k)}$. This is in the compact hull since the unused weight can be placed at 0. Define $T_k(s)=1$ if the $i(k)$th component of $\sum_{l\ne k} L_l s_l p_{i(l)}$ is at least $a_k$, and $0$ otherwise. All entries of the $p_i$ are nonnegative, so $T$ is order preserving. Iterate from the all-zero vector to obtain a nondecreasing sequence and hence a fixed point. There the $k$th tested coordinate is either $<a_k$ or $\ge a_k+L_k>b_k$. Thus $x(s)$ avoids every interval.

Apply the AHKM rounding argument \cite[Thm.~3.1]{AHKM}. To see it directly take $U$ uniform on $[0,1]$, $V=(U+1/2)\bmod 1$, and $W=3/2-U-V$. The two half intervals each map linearly onto $[0,1]$ under the rule for $W$, so $W$ is uniform too. The triple $Y=\tfrac23(U,V,W)$ sums to one, with uniform marginals on $[0,2/3]$. For any fractional cover, $\{t:f_i(t)\ge Y_i\}$ form an ordinary cover. Averaging gives $1\le(3/2)\sum_i\int\min(f_i,2/3)\le(3/2)\sum_i\int f_i$.
\end{proof}
\begin{proof}[Proof of Theorem~\ref{main}] Assign a median vertex in each coordinate, ties arbitrarily. If the three assignments differ, integrate the cover against the law from Lemma~\ref{distinct}. If not, an unassigned vertex is a cube vertex: in a packed triangle any strictly intermediate coordinate value would be the unique median and would have been assigned. Apply Lemma~\ref{corner}. On the standard simplex $\sum x_i=1$, the functions $f_i(x)=(2/3-x)_+$ form a cover, with cost $2/3$.\end{proof}
\section*{A separate cyclic gluing criterion}
Suppose a convex cyclically invariant $K\subseteq[0,1]^3$ contains $p=(0,a,b)$ and $q=(1,c,e)$. Put $h_0=(a+b)/3$, $h_1=(1+c+e)/3$, $D=h_1-h_0$. If $D\ge1/3$, there is a measure on $K$ of mass $\ge2/3$ whose marginals are dominated by Lebesgue measure, even for full-dimensional $K$.

Here is a self-contained orbit calculation. Any orbit triangle whose vertex values range over $0,s,W$ after translation or reflection has coordinate center $h$ and supports a law uniform on each $[h-W/3,h+W/3]$. Indeed its cyclic triangular hull is described in barycentric coordinates by the cyclic permutations of $(0,s,W)$. For the target vector $(h,h+r,h-r)$ the barycentric weights equal $1/3$ plus $r$ times a permutation of $(s,W-s,-W)/(s^2-sW+W^2)$, up to a common sign (by solving the circulant linear system). Choose the sign so that the largest magnitude of a negative slope is $\max(s,W-s)/(s^2-sW+W^2)$. This permits $r\ge W/3$, since $s^2-sW+W^2\ge W\max(s,W-s)$. Its orbit contains a triangle whose equal edge masses give uniform marginals on $[h-W/3,h+W/3]$, after contracting to that radius. If $W=0$ use zero mass below.

Thus the two orbit triangles yield uniform radii $R_0=\max(a,b)/3$ and $R_1=\max(1-c,1-e)/3$, respectively. Contract about the respective centers to radii $r_i\le R_i$ with $r_0+r_1=\min(D,R_0+R_1)$ and give each law mass $2r_i$. Their intervals are interior-disjoint. The diagonal between the centers lies in $K$; put length measure on it just across the remaining gap. All marginals now equal Lebesgue measure restricted to the joined interval of length $D+\min(D,R_0+R_1)$. Since $R_0\ge h_0/2$ and $R_1\ge(1-h_1)/2$, the length is at least $\min(2D,(1+D)/2)\ge2/3$. No assertion for general cyclic bodies with $D<1/3$ follows.

\paragraph{Scope and novelty.} Earlier project notes supplied the signed-chord criterion used in Lemma~\ref{corner}, a cyclic planar result, and the uniform edge law. The distinct-median clipping construction, its assembly into the unrestricted triangle theorem, and the full-dimensional gluing criterion are the contributions of this note. The general three-dimensional conjecture remains unresolved here. An independent priority claim would require further review.
\begin{thebibliography}{9}
\bibitem{AHKM} R. Aharoni, R. Holzman, M. Krivelevich and R. Meshulam, Fractional planks, \emph{Discrete Comput. Geom.} 27 (2002), 585--602. \href{https://doi.org/10.1007/s00454-001-0088-x}{DOI}.
\bibitem{BMP} P. Brass, W. Moser and J. Pach, \emph{Research Problems in Discrete Geometry}, Springer, 2005. \href{https://doi.org/10.1007/0-387-29929-7}{DOI}.
\end{thebibliography}
\end{document}
